% TOP_PROBLEM: {"id": 30006726, "title": "Green-Griffiths-Lang Conjecture on entire curves in varieties of general type", "category_id": 5, "category": {"id": 5, "name": "algebraic_geometry", "display_name": "Algebraic Geometry", "description": "Geometric objects defined by polynomial equations.", "slug": "algebraic-geometry", "order_index": 5, "created_at": "2026-07-31T15:26:25.671Z"}, "status": "open"}
% FIRST_POSTED: 2026-09-25
% ENTRY_KIND: statement_only
% !TeX program = lualatex
% Definition and sources only; this file is not an LLM solution attempt.
\documentclass[11pt]{article}
\usepackage[margin=25mm]{geometry}
\usepackage{fontspec}
\setmainfont{Latin Modern Roman}
\usepackage{amsmath,amssymb,mathtools}
\usepackage{unicode-math}
\setmathfont{Latin Modern Math}
\usepackage{xurl,hyperref}
\hypersetup{colorlinks=true,urlcolor=blue}
\setcounter{secnumdepth}{0}
\setlength{\parindent}{0pt}
\setlength{\parskip}{5pt}
\setlength{\emergencystretch}{3em}
\begin{document}
\section*{\texorpdfstring{Green-Griffiths-Lang Conjecture on entire curves in varieties of general type}{Green-Griffiths-Lang Conjecture on entire curves in varieties of general type}}\label{30006726}
\subsection{Definitions and mathematical statement}
An entire curve is a nonconstant holomorphic map $f:{\mathbb{C}}\to X$. Algebraic degeneracy means $\overline{f({\mathbb{C}})}^{\rm Zar}\subsetneq X$. The catalogue's literal assertion quantifies a possibly different proper subvariety for each curve:
\[
 \forall f:{\mathbb{C}}\to X\text{ entire},\quad\exists Z_f\subsetneq X:\ f({\mathbb{C}})\subseteq Z_f.
\]
The stronger usual uniform Green--Griffiths--Lang formulation places \emph{all} such curves in one proper algebraic subset depending only on $X$; Demailly's Introduction states that formulation. The existence of jet-differential equations is an important method but not by itself equivalent to either degeneracy conclusion: one must control their common solution locus. These distinctions are recorded rather than silently repairing the source wording.
\par\smallskip\textit{Formulation and concept references:} \href{https://arxiv.org/abs/1011.3636}{[S1]}, \href{https://arxiv.org/html/1011.3636v2}{[E1]}.

\subsection{Short English statement}
Must every entire curve in a projective variety of general type lie inside a proper algebraic subset, with the stronger standard version asking for one subset covering them all?
\subsection{Sources}
\begin{itemize}
\item[S1]J.-P. Demailly, Holomorphic Morse inequalities and the Green-Griffiths-Lang conjecture, arXiv:1011.3636 [math.AG], 2010.
\url{https://arxiv.org/abs/1011.3636}.
\textit{Formulation source.}
\item[S2]Computing Jet Differentials and the Green-Griffiths-Lang Conjecture for Complements of Smooth Plane Curves.
\url{https://arxiv.org/abs/2608.20479}.
\textit{Further reference; not a proof endorsement.}
\item[S3]Non-reductive geometric invariant theory and hyperbolicity.
\url{https://link.springer.com/article/10.1007/s00222-023-01219-z}.
\textit{Further reference; not a proof endorsement.}
\item[E1]Jean-Pierre Demailly, Holomorphic Morse inequalities and the Green-Griffiths-Lang conjecture, arXiv:1011.3636v2 (2010), Introduction, Conjecture (0.2) and discussion of jet differentials.
\url{https://arxiv.org/html/1011.3636v2}.
\textit{Added primary source: Quantifier distinction and limitations of the jet-differential reformulation. Consulted 24 September 2026.}
\end{itemize}
\end{document}
